\chapter{Cefalexin (Cephalexin)}

\section*{Definition and Overview}
Cefalexin, also spelled cephalexin, is an oral antibiotic belonging to the first-generation cephalosporin family. It is used to treat a range of common bacterial infections, including infections of the skin, soft tissues, bones, urinary tract, and respiratory tract. Like other cephalosporins, it works by interfering with the bacterial cell wall, causing susceptible bacteria to die.

Because it is absorbed well from the gut and has a relatively narrow spectrum, cefalexin is often used as a first-line oral agent for straightforward community-acquired infections. It is available only on prescription, and the choice of antibiotic should be guided by a doctor or pharmacist.

\section*{Epidemiology}
Antibiotics such as cefalexin are among the most commonly prescribed medicines worldwide. Skin infections caused by \textit{Staphylococcus} and \textit{Streptococcus} bacteria, urinary tract infections, and middle-ear infections account for a large share of prescriptions, particularly in children and older adults.

Prescribing patterns vary between countries and are shaped by local rates of antibiotic resistance. Overuse of broad-spectrum antibiotics has led to rising resistance among common bacteria, which is why cefalexin -- with its relatively targeted spectrum -- remains a useful option for many mild to moderate infections.

\section*{Aetiology and Pathophysiology}
Cefalexin is used when an infection is caused by bacteria that are susceptible to first-generation cephalosporins:

\begin{itemize}
    \item \textbf{Skin and soft tissue infections:} usually due to \textit{Staphylococcus aureus} and \textit{Streptococcus pyogenes}
    \item \textbf{Urinary tract infections:} often caused by \textit{Escherichia coli} and other Gram-negative bowel organisms
    \item \textbf{Respiratory tract infections:} such as tonsillitis and sinusitis, often streptococcal in origin
    \item \textbf{Bone and joint infections:} where oral step-down therapy is appropriate
\end{itemize}

Like other beta-lactam antibiotics, cefalexin binds to enzymes that build the bacterial cell wall, weakening the wall and leading to destruction of the bacterium. It does not work against viruses, fungi, or bacteria that produce enzymes able to break the drug down.

\section*{Clinical Features}
Symptoms of the infections treated with cefalexin vary with the site involved:

\begin{itemize}
    \item \textbf{Skin infections:} redness, warmth, swelling, tenderness, and discharge from an infected wound or boil
    \item \textbf{Urinary tract infections:} burning or pain on passing urine, frequent urination, urgency, and cloudy or smelly urine
    \item \textbf{Throat and tonsil infections:} sore throat, fever, swollen neck glands, and difficulty swallowing
    \item \textbf{Strep throat in particular:} sudden severe sore throat with fever, headache, and a lack of cough
    \item \textbf{General:} fever, malaise, and feeling generally unwell
\end{itemize}

Side effects of the drug itself may include nausea, diarrhoea, abdominal discomfort, and rash, which can be difficult to tell apart from the underlying illness.

\section*{Differential Diagnosis}
The decision to use cefalexin rests on identifying whether a bacterial infection is present and which antibiotic is appropriate:

\begin{itemize}
    \item \textbf{Viral infections:} most sore throats, sinus infections, and common colds are viral and will not respond to antibiotics
    \item \textbf{Fungal infections:} skin rashes may be fungal (such as tinea) rather than bacterial
    \item \textbf{Non-infectious causes:} drug reactions, eczema, or dermatitis may mimic bacterial skin infection
    \item \textbf{Resistant bacteria:} infections due to methicillin-resistant \textit{Staphylococcus aureus} (MRSA) or penicillin-resistant organisms need other agents
    \item \textbf{Deep infections:} abscesses or infections under the skin surface may require drainage rather than antibiotics alone
\end{itemize}

\section*{When to Seek Medical Care}
See a doctor if an infection is severe, spreading rapidly, or not improving within a couple of days of starting treatment, or if you have side effects that concern you.

\textbf{Call 000 (or local emergency services) for:}
\begin{itemize}
    \item Difficulty breathing, wheezing, or swelling of the face, lips, or throat (possible severe allergic reaction)
    \item A widespread rash, especially with fever or peeling skin
    \item Signs of sepsis: high fever, confusion, rapid breathing, or mottled skin
    \item Severe bloody diarrhoea, which may indicate an antibiotic-related bowel infection
\end{itemize}

\section*{Diagnosis and Investigations}
\begin{itemize}
    \item \textbf{Clinical assessment:} history and examination to identify the site and severity of infection
    \item \textbf{Wound or throat swabs:} culture to identify the organism and its antibiotic sensitivities
    \item \textbf{Urine testing:} dipstick and urine culture for suspected urinary tract infection
    \item \textbf{Blood tests:} full blood count and inflammatory markers (CRP) in more severe or systemic infection
    \item \textbf{Imaging:} X-ray, ultrasound, or MRI if a deep collection or bone infection is suspected
    \item \textbf{Allergy review:} careful history of previous reactions to penicillins or cephalosporins before prescribing
\end{itemize}

\section*{Management}
\begin{itemize}
    \item \textbf{Antibiotic course:} cefalexin is usually taken by mouth several times a day; the dose, interval, and duration depend on the infection, age, and kidney function
    \item \textbf{Complete the course:} finish the full prescribed course even if symptoms improve, to clear the infection and reduce resistance
    \item \textbf{Pain and fever control:} paracetamol or ibuprofen can help with discomfort while the antibiotic takes effect
    \item \textbf{Local care:} clean and dress infected skin wounds; some abscesses need drainage by a clinician
    \item \textbf{Fluid and rest:} adequate fluids and rest support recovery
    \item \textbf{Follow-up:} repeat review if symptoms persist, or if cultures show a resistant organism
\end{itemize}

\section*{Complications}
\begin{itemize}
    \item Allergic reactions ranging from a mild rash to rare, life-threatening anaphylaxis
    \item Gastrointestinal upset, including nausea and diarrhoea
    \item \textit{Clostridium difficile} infection: antibiotic-related bowel inflammation causing severe diarrhoea
    \item Yeast infections (such as oral or vaginal thrush) due to changes in normal flora
    \item Antibiotic resistance developing in the community with inappropriate use
    \item Failure of treatment if the causative organism is resistant
\end{itemize}

\section*{Prognosis and Outcomes}
Most common bacterial infections treated with cefalexin resolve within days of starting the medicine, with improvement often seen within 48 hours. Outcomes are generally excellent when the organism is susceptible and the infection is not deep or complicated. Prognosis is less favourable in people with weakened immunity, poor circulation, or infections involving bone or deep tissue, and in those with resistant organisms. Side effects usually settle once the course is finished.

\section*{Prevention and Self-Care}
\begin{itemize}
    \item Use antibiotics only when prescribed and exactly as directed
    \item Do not save leftover antibiotics or share them with others
    \item Practise good hand hygiene and keep wounds clean and covered to prevent infection
    \item Complete prescribed courses fully to discourage resistance
    \item Maintain up-to-date vaccinations, which reduce the risk of vaccine-preventable bacterial illness
    \item Tell your doctor about any past allergic reaction to antibiotics before taking a new one
    \item Seek early review if a skin or urinary infection is not settling with treatment
\end{itemize}

\section*{Sources and Further Reading}
The following organisations provide reliable information on cefalexin and the rational use of antibiotics.

\begin{itemize}
    \item National Health Service. \textit{NHS conditions guide on antibiotics and their uses.} https://www.nhs.uk/conditions/
    \item Mayo Clinic. \textit{Antibiotics: safe use, prescribing, and side effects.} https://www.mayoclinic.org/
    \item Centers for Disease Control and Prevention. \textit{Antibiotic prescribing and use information.} https://www.cdc.gov/
    \item World Health Organization. \textit{Antimicrobial resistance and rational antibiotic use.} https://www.who.int/
    \item MedlinePlus. \textit{Consumer health information on cefalexin (cephalexin).} https://medlineplus.gov/
    \item MSD Manual. \textit{Cephalosporins and beta-lactam antibiotics -- professional reference.} https://www.msdmanuals.com/
\end{itemize}
